\documentclass[book.tex]{subfiles}

\begin{document}


\chapter{Ising model}
\label{chap:ising_model}
\noindent\rule{11cm}{0.4pt} \\
\textbf{Prerequisites:}   \autoref{chap:stat_thermo}, \autoref{chap:noninteracting}  \\
\textbf{Difficulty Level:} ** \\
\noindent\rule{11cm}{0.4pt}
\vspace{5mm}

Thus far, our analyses have focused on non-interacting systems. Non-interacting systems require the evaluation of a single-molecule partition function, which generally is tractable (either exactly or approximately). However, very few practical problems involve either non-interacting molecules or approximately non-interacting molecules. Furthermore, a variety of important physical phenomena are not exhibited in non-interacting systems. For example, phase transitions are a hallmark issue in interacting systems. Our goal for this chapter is to discuss the effect of interactions on thermodynamic behavior.

We introduce the Ising Model, which constitutes the basis for the analysis of long ranged correlations. Behind its apparently simple formulation, the Ising model hides a complex mathematical problem, specially for 2 or 3 dimensional systems, that requires a number of approximations in order to derive a reasonable descriptive model for the study of these correlations. Of course, there is always a cost associated for using these approximations, and we will discuss these limitations over the course of this chapter. For now, the behavior of the Ising model in 1D system will be presented, and in the second part, we will introduce the mean field approximation for the study of 2D systems. 

\begin{marginfigure}
% \includegraphics[width=\linewidth]{figures/Physics/statistical_mechanics/ising_model/lattice.png}
\includegraphics[width=\linewidth]{figures/Physics/statistical_mechanics/ising_model/2D_lattice_structure.png}
\caption{Schematic representation of different 2D lattice structures}
\label{lattice}
\end{marginfigure}

Consider a lattice of $N$ sites. The definition of the lattice can include many different geometries (square lattice, triangular lattice, etc..) and dimensionality (1 dimension, 2 dimensions, etc...) (Fig.~\ref{lattice}). In the presence of a magnetic field, $h$, the energy of the system in a particular state, $\nu$, is:
\begin{align}
E_\nu &= -\sum_{i=1}^N h s_i + (\text{energy due to interactions between spins})
\end{align}
According to the expression, each site contains exactly 1 spin associated to its magnetic moment $\pm \mu$. Thus, the net magnetic moment is defined by the spin state $s_i$, that is either spin up ($s_i = +1$) or spin down ($s_i = -1$).

A simple model for the interaction energy is:
\begin{align}
- J \sum_{\langle i j \rangle}s_i s_j
\end{align}

\noindent where the sum over $\langle i j \rangle$ implies a summation over all nearest neighbor pairs in the lattice, with a coupling strength $J$. 

The system energy is now defined to be:
\begin{align}\label{Eq1}
E_\nu= -h\sum_{i=1}^N s_i - J \sum_{\langle i j \rangle}s_i s_j
\end{align}

\noindent which is the Ising Model. In general, this model can be applied to a variety of problems (magnetic systems, binary alloys, liquid-gas phase transition, neural networks - Hopfield model).

Notice that when $J> 0$, it is energetically favorable that the neighboring spins remain aligned. At low temperatures, this alignment between spins will lead to a cooperative phenomenon called spontaneous magnetization. 

\begin{marginfigure}
% \includegraphics[width=\linewidth]{figures/Physics/statistical_mechanics/ising_model/ferromagnetic_ground_states.png}
\includegraphics[width=\linewidth]{figures/Physics/statistical_mechanics/ising_model/ground states_2D_ising_model.png}
\caption{Ground states (minimum energy) for the 2-dimensional
Ising model (square lattice) for a ferromagnetic system (with coupling constant $J >0$) and an antiferromagnetic system (with coupling constant $J < 0$).}
\label{fig1}
\end{marginfigure}
In general the sign of the coupling constant $J$ dictates whether the spins prefer to align with their neighbors or anti-align (Fig.~\ref{fig1}):

\begin{itemize}
\item Ferromagnetic: the spins prefer to align with their neighbors, occurring when $J > 0$.
\item Antiferromagnetic: the spins prefer to anti-align with their neighbors, occurring when $J < 0$.
\end{itemize}

Define the total magnetization $M$ of the system to be:
\begin{align}\label{Eq2}
M = \bigg\langle \sum_{i=1}^N s_i \bigg\rangle = \frac{\partial \log Q}{\partial \beta h}
\end{align}

\begin{marginfigure}
% \includegraphics[width=\linewidth]{figures/Physics/statistical_mechanics/ising_model/spontaneous_magnetization.png}
\includegraphics[width=\linewidth]{figures/Physics/statistical_mechanics/ising_model/spont_magnetization.png}
\caption{Spontaneous Magnetization}
\label{fig2a}
\end{marginfigure}

In the absence of an external magnetic field $h$, for $J>0$ we expect the system to develop a net magnetization $\langle M \rangle$ at temperatures below $T_c$, defined as the Curie Temperature or Critical Temperature (Fig.\ref{fig2a})

\subsection*{\textit{Non-interacting case ($J = 0$)}}

The partition function for the Ising model with $J = 0$ is given by
\begin{align}
Z &= \sum_{s_1=-1,1}\sum_{s_2=-1,1}...\sum_{s_N=-1,1} \exp\left(\beta h\sum_{i=1}^Ns_i\right)\\
&=\prod_{j=1}^N \left(\sum_{s_j=-1,1}e^{\beta h s_j}\right) =\left(\sum_{s=-1,1}e^{\beta h s}\right)^N \\
&= \left(e^{\beta h} + e^{-\beta h}\right)^N \\
&= 2^N \cosh^N\left(\beta h \right) 
\end{align}

\noindent which is valid for any lattice type and dimensionality. The total magnetization $M$ is:
\begin{align}\label{Eq2}
M = \frac{\partial \log Z}{\partial \beta h} = N\tanh \left(\beta h \right) 
\end{align}

The average energy $\langle E \rangle$ and entropy $S$ are given by:
\begin{align}
\langle E \rangle &= -h \bigg\langle \sum_{i=1}^Ns_i \bigg\rangle \\
&= -hM \\
&= -hN\tanh \left(\beta h \right) \\
S &= \frac{E-F}{T} \\
&= \frac{-hN\tanh \left(\beta h \right) + Nk_BT\log\left[2\cosh \left(\beta h \right)\right]}{T} 
\end{align}

Note, $\langle E \rangle \rightarrow -hN$ and $S\rightarrow 0$ as $T \rightarrow 0$, and $\langle E \rangle \rightarrow 0$ and $S \rightarrow Nk_B \log 2$ as $T \rightarrow \infty$.

\begin{marginfigure}
\includegraphics[width=\linewidth]{figures/Physics/statistical_mechanics/ising_model/total_magnetization.png}
\caption{The total magnetization $M=N$ versus the external field $\beta h$ for the non-interacting Ising model $J = 0$).}
\label{fig2}
\end{marginfigure}

The non-interacting Ising model ($J =0$) exhibits a gradual change
in the magnetization $M$ with the external field $h$ (Fig.~\ref{fig2}). Thus, no phase transition is exhibited in the non-interacting case; there is never a phase transition for non-interacting molecules. We now turn to cases where the interactions are turned on, such that $J \neq 0$.

\section*{\textit{Ising model in 1 dimension}}

Consider a line of spins in $1$ dimension with periodic boundary
conditions such that $s_N$ interacts with $s_1$. For the Ising model in $1$ dimension, the energy is written as:
\begin{align}\label{Eq3}
\beta E = -\tilde{h} \sum_{i=1}^N s_i -\tilde{J}\sum_{i=1}^N s_is_{i+1}
\end{align}
\noindent where 
\begin{align}
\tilde{h} & =\beta h  \\
\tilde{J} &= \beta J
\end{align}
Define a transfer function $T(s_i,s_{i+1})$ to be
\begin{align}\label{Eq4}
T(s_i,s_{i+1}) &= \exp \left[ \tilde{h}\left(\frac{s_i}{2}+\frac{s_{i+1}}{2}\right)+\tilde{J}s_is_{i+1}\right]
\end{align}

With this definition, we write the partition function as:
\begin{align}
Z = \sum_{s_1=-1,1}\sum_{s_2=-1,1}...\sum_{s_N=-1,1} T(s_1,s_2)T(s_2,s_3)...T(s_N,s_1)
\end{align}

Now construct the \textit{transfer matrix} to be:
\begin{align}\label{Eq5}
T = \left[ \begin{array}{cc} T(1,1) & T(1,-1) \\ T(-1,1) & T(-1,-1) \end{array} \right] = \left[ \begin{array}{cc} e^{\tilde{h}+\tilde{J}} & e^{-\tilde{J}} \\ e^{-\tilde{J}} & e^{-\tilde{h}+\tilde{J}} \end{array} \right] 
\end{align}

The partition function is then written as:
\begin{align}\label{Eq6}
Z = \sum_{s_1=-1,1}T^N(s_1,s_1) = Tr\{T^N\}
\end{align}

\noindent where $Tr\{T^N\}$ indicates a trace of the matrix $T^N$.

The trace of a matrix is the sum of its eigenvalues, and the eigenvalues of $T^N$ are $\lambda_\pm ^N$, where $\lambda_\pm$ are the eigenvalues of $T$. In our case:
\begin{align}\label{Eq7}
\lambda_\pm =e^{\tilde{J}} \cosh \tilde{h} \pm \left(e^{-2\tilde{J}} + e^{2\tilde{J}}\sinh^2\tilde{h}\right)^{1/2}
\end{align} 

Therefore, the exact solution for the partition function is:
\begin{align}
Z = \left[e^{\tilde{J}} \cosh \tilde{h} + \left(e^{-2\tilde{J}} + e^{2\tilde{J}}\sinh^2\tilde{h}\right)^{1/2}\right]^N + \\ \left[e^{\tilde{J}} \cosh \tilde{h} \pm \left(e^{-2\tilde{J}} + e^{2\tilde{J}}\sinh^2\tilde{h}\right)^{1/2}\right]^N 
\end{align}

\begin{marginfigure}[-5cm]
% \includegraphics[width=\linewidth]{figures/Physics/statistical_mechanics/ising_model/phase_transition_1d.png}
\includegraphics[width=\linewidth]{figures/Physics/statistical_mechanics/ising_model/Phase_transition_1D (1).png}
\label{fig3a}
\end{marginfigure}
\begin{marginfigure}
% \includegraphics[width=\linewidth]{figures/Physics/statistical_mechanics/ising_model/phase_transition_2d.png}

\includegraphics[width=\linewidth]{figures/Physics/statistical_mechanics/ising_model/Phase_transition_2D (1).png}
\caption{Energy cost of a phase transition from an ordered phase to
a disordered phase in 1D and 2D.}
\label{fig3b}
\end{marginfigure}

For $\tilde{J} =0$, the partition function $Z=2^N \cosh^N \tilde{h}$, which agrees with our previous result. For $\tilde{h}=0$, the partition function $Z=2^N\cosh^N\tilde{J} + 2^N\sinh^N\tilde{J} \approx 2^N\cosh^N\tilde{J}$ for large $N$.

Since $\lambda_+ > \lambda_-$, the partition function is $Z \approx \lambda_+^N$ for $N \gg 1$. The average magnetization is given by 
\begin{align}
M = \frac{\partial \log Z}{\partial h} = N \frac{\sinh \tilde{h}}{\sqrt{e^{-4\tilde{J}}}+\sinh^2\tilde{h}}
\end{align}
In the limit $\tilde{h} \rightarrow 0^+$, the magnetization approaches zero. Thus, this exact solution for the 1D Ising model does not exhibit a phase transition from a disordered state to an ordered state. The physical justification for this lies in the fact that forming a disordered phase from an ordered phase occurs through an energy change $\Delta E = 4J$ in 1D (independent of $N$). Thus, the net magnetization per particle should vanish for all temperatures higher than $T\sim J/Nk_B$, which is very small considering that $N$ is large for macroscopic systems (Fig.~\ref{fig3b}a). 

In higher dimensions ($D \geq 2$), the energy cost scales with $N$ (\textit{e.g.} $\Delta E \sim N^{1/2}J$ in 2D), thus leading to a finite-$T$ phase transition (Fig.~\ref{fig3b}b). 
\newpage

\subsection*{Mean-field approximation}
The solution for the Ising model for $D \geq 2$ is not as easily found. In 1944, Lars Onsager found an exact solution for the 2D
Ising model for $h = 0$ that exhibits a spontaneous magnetization
for sufficiently large $J$. However, there is no exact solution for the Ising model for $D \geq 3$. 

We introduce the mean-field approximation. The idea of this method is to focus on one particular particle of the system, and assume that the role of the neighboring particles is to form an average molecular magnetic field, which acts on the tagged particle (Fig.~\ref{fig4}). In the next lecture, we test the limits of this approximation.

\begin{marginfigure}[-6cm]
% \includegraphics[width=\linewidth]{figures/Physics/statistical_mechanics/ising_model/mean_field_approx.png}
\includegraphics[width=\linewidth]{figures/Physics/statistical_mechanics/ising_model/Mean field.png}
\caption{Schematic of the mean-field approximation.}
\label{fig4}
\end{marginfigure}

Define the average magnetization per site 
\begin{align}
m=\frac{M}{m}=\frac{1}{N}\sum_{i=1}^N\langle s_i\rangle = \langle s\rangle
\end{align}
by noting that spatial invariance implies $\langle s_i\rangle = \langle s\rangle$. Define the coordination number z as the number of nearest neighbors
to any given site divided by $2$ (\textit{e.g.} $z = 2$ for 2D square lattice and $z = 3$ for 3D square lattice). The mean-field approximation replaces the near-neighbor interactions
by the average interaction, thus $s_j \approx \langle s_j\rangle$, giving the energy approximation:

\begin{marginfigure}
\includegraphics[width=\linewidth]{figures/Physics/statistical_mechanics/ising_model/tanh_K_plot.png}
\caption{Plots of $\tanh (Km)$ versus $m$ for various values of $K$. Circles mark points that satisfy the self-consistent equation $m = \tanh (Km)$.}
\label{fig5}
\end{marginfigure}
\begin{align}\label{Eq8}
E &= -h\sum_{i=1}^N s_i - J\sum_{\langle i j\rangle}s_is_j \\
&\approx -h\sum_{i=1}^N s_i - Jzm\sum_{i=1}^N s_i
\end{align}

\noindent which reduces the problem to an effective non-interacting model.

Therefore, we write the average magnetization per site as:
\begin{align}\label{Eq9}
m = \frac{\sum_{s=-1,1}s \exp\left[ (\beta h + \beta Jzm) s \right]}{\sum_{s=-1,1} \exp\left[ (\beta h + \beta Jzm) s \right]} = \tanh (\beta h + \beta Jzm)
\end{align}

\noindent which performs the average of the site magnetization using the mean-field site energy in the Boltzmann weight for the spin probability.

For the case $h = 0$ (no external field), the magnetization satisfies
\begin{align}
m &= \tanh (Km)
\end{align}
where $K = \beta Jz$. This is a self-consistent mean-field equation. Local order is dictated by the surrounding order that is itself defined by the local order, leading to a self-consistent condition that must be met. Two scenarios exist (Fig.~\ref{fig5}):

\begin{marginfigure}
% \includegraphics[width=\linewidth]{figures/Physics/statistical_mechanics/ising_model/ising_phase_diagram.png}
\includegraphics[width=\linewidth]{figures/Physics/statistical_mechanics/ising_model/ising_phase_diagram0.png}
\caption{Phase diagram for the mean-field solution for the Ising
model, showing $m$ versus $T/T_c$.}
\label{fig6}
\end{marginfigure}

\begin{itemize}
\item For $K < 1$, there is only one solution for m that corresponds
to $m = 0$.
\item For $K >1$, there are 3 solutions for $m$ corresponding to $m=0 $ and $m=\pm m^*$. The solution $m=0$ is unphysical as it is a free
energy maximum.
\end{itemize}

Therefore, $K = K_c = 1$ defines a critical temperature $T_c = Jz/k_B$, and we can write $K = T_c/T$ (Fig.~\ref{fig6}). 

Near $T = T_c$, we can expand the self-consistent equation near $m = 0$ to get:
\begin{align}
m = \tanh \left( \frac{T_c}{T}m \right) \approx \frac{T_c}{T}m - \frac{1}{3}\left( \frac{T_c}{T}\right)^3m^3 \rightarrow m \approx \pm \sqrt{\frac{3}{T_c}}(T_c-T)^{1/2} 
\end{align}

Generally, the external field $h$ acts as the intensive variable that
is conjugate to the magnetization $M$ (extensive variable). Earlier in the book, we learned that thermodynamic stability dictates that a derivative of an intensive variable with respect to its conjugate
extensive variable is a positive quantity, \textit{e.g.}:
\begin{align}\label{Eq10}
\left( \frac{\partial T}{\partial S} \right)_{V,N} & > 0\\ 
- \left( \frac{\partial p}{\partial V} \right)_{T,N} &> 0\\
\left( \frac{\partial \mu_i}{\partial N_i} \right)_{T,V,N_{j\neq i}} &> 0
\end{align}

Similarly, thermodynamic stability dictates that:
\begin{align}\label{Eq11}
\left( \frac{\partial \beta h}{\partial M} \right)_{\beta Jz,N}>0
\end{align}

\begin{marginfigure}
\includegraphics[width=\linewidth]{figures/Physics/statistical_mechanics/ising_model/ising_plot_m_bh.png}
\caption{Plot of $m$ versus $\beta h$ for several values of $K = \beta Jz = T_c=T$. The dashed part of each curve indicates where the system is unstable.}
\label{fig7}
\end{marginfigure}

The phase diagram for the mean-field solution is dictated by the
self-consistent mean field equation (Fig.~\ref{fig7}):
\begin{align}\label{Eq12}
m = \tanh (\beta h + \beta Jzm)
\end{align}

This equation of state can be likened to the vapor-liquid equation
of state $p = p(v)$.

The phase diagram for the mean-field approximation for the Ising model exhibits similar characteristics as the liquid-vapor phase diagram. The correspondence between $M-h$ in the magnetic system and $p-V$ in the liquid-vapor system reflects the universal physical issues that are addressed with the Ising model and why it can be used to understand a variety of phenomena (Fig.~\ref{fig9}).

\begin{marginfigure}
\includegraphics[width=\linewidth]{figures/Physics/statistical_mechanics/ising_model/ising_hysteresis.png}
\caption{Plot of $m$ versus $\beta h$ for $K = \beta Jz = T_c/T = 1.5$. The arrows indicate the hysteresis in magnetization.}
\label{fig9}
\end{marginfigure}

\end{document}